\subsection{Наблюдаемые результаты разработки}
Заменяющая матрица завершила все 40 назначенных генераций и 72 хода CLI. Аудит сверяет записанные запросы и полную промежуточную историю с зафиксированным манифестом и подтверждает 387\,125 входных и выходных токенов. Их суммарная оценка по опубликованному тарифу составляет 0.7417401 USD, чувствительность без кеша --- 0.8381625 USD. В принятой матрице не обнаружены события инструментов, повторов или сжатия истории. Это реальные сгенерированные решения, отличные от программных фикстур и невыполненных подтверждающих манифестов.

\begin{table}[htbp]
\centering\small
\caption{Наблюдаемые результаты разработки: восемь назначенных задач на режим. Стоимость --- средняя условная цена решения по API, токены включают накладные расходы CLI. Диапазон качества отражает неразрешённые исходы, а не доверительный интервал.}
\label{tab:codexpilot}
\begin{tabular}{lrrrr}
\toprule
Режим & Токены & USD & Успех/оценено & Диапазон \\
\midrule
1 / нейтр. & 5884 & 0.01413 & 4/7 & 0.500--0.625 \\
1 / роли & 5437 & 0.01196 & 3/7 & 0.375--0.500 \\
3 / нейтр. & 16304 & 0.02863 & 4/7 & 0.500--0.625 \\
3 / роли & 16041 & 0.02871 & 2/7 & 0.250--0.375 \\
Прямой & 4725 & 0.00929 & 4/7 & 0.500--0.625 \\
\bottomrule
\end{tabular}
\end{table}


\begin{figure}[htbp]
\centering\includegraphics[width=\textwidth]{figures/codex_pilot.pdf}
\caption{Все кандидаты завершённой development-матрицы. Точки показывают восемь назначенных задач в каждом режиме, горизонтальные отметки --- средние. Горизонтальное смещение лишь разделяет точки и не является интервалом неопределённости. Обе панели включают задачу с недоступной тестовой оценкой. Слева --- тысячи токенов, справа --- условная API-цена в центах США.}
\label{fig:codexobservations}
\end{figure}

В каждом режиме семь задач пригодны к оценке, одна недоступна. Эталон /1005 требует сети и не проходит в среде с отключённой сетью; результат модели здесь недоступен во всех режимах, хотя исходные тесты завершаются ошибкой. Остальные семь эталонов проходят, а все восемь заведомо неправильных контролей проваливаются. Среди 35 оцениваемых кандидатов 17 успешны: прежний потолок качества в этой development-выборке отсутствует. Однако это семь оцениваемых кластеров задач, а не 35 независимых задач.

Прямой решатель имеет минимальную наблюдаемую среднюю условную цену (0.00929 USD) и проходит 4/7 оцениваемых задач, как и оба нейтральных протокола. Одно обращение с ролями даёт 3/7, три обращения с ролями --- 2/7. Относительно соответствующего нейтрального режима роли уменьшают среднее число токенов на 7.6\% при одном обращении и на 1.6\% при трёх. Три обращения расходуют в 2.77 раза больше токенов, чем одно, при нейтральных метках, и в 2.95 раза больше при ролях. Эти описательные наблюдения не устанавливают неухудшение, статистически доказанный вред ролей или ранжирование на всём бенчмарке.

\begin{table}[htbp]
\centering\small
\caption{Исход каждой назначенной задачи. P --- успех, F --- ошибка, T --- тайм-аут, U --- оценка недоступна после проверки эталона и отрицательного контроля. Задачи не удаляются.}
\label{tab:codextasks}
\begin{tabular}{lccccc}
\toprule
Задача & 1 / нейтр. & 1 / роли & 3 / нейтр. & 3 / роли & Прямой \\
\midrule
325 & P & P & P & F & P \\
322 & F & F & F & F & F \\
1036 & F & F & F & F & F \\
1005 & U & U & U & U & U \\
361 & P & P & P & P & P \\
1087 & P & F & P & P & P \\
45 & F & F & F & F & F \\
309 & P & P & P & F & P \\
\bottomrule
\end{tabular}
\end{table}


Один ответ одновызовного ролевого режима на /1087 нарушает заранее установленный формат извлечения и учитывается как неуспех; более выгодный блок кода вручную не выбирается. Среди ошибок трёхвызовного ролевого режима тесты /325 указывают на удаление скобок, которые следовало сохранить, а /309 --- на несовпадение оси нормализации. Это проверка поведения финального кандидата, а не доказательство независимости внутренних ролей или того, что ошибку вызвал именно этап ревью. Полные выведенные тексты этапов и трассы утверждений доступны для дальнейшей явно поисковой диагностики.

Проверка общих ошибок после оценки выявила ограничения валидности. Тест /1036 требует буквального заголовка графика, отсутствующего в предоставленном условии; /45 содержит противоречивые требования к именам столбцов и осей, а тесты выбирают один вариант; /322 подменяет конкретные функции работы с файлами и процессами и тем самым штрафует альтернативы, явно не запрещённые условием. Успех эталона таких несоответствий не обнаруживает. Исходные тестовые оценки и все назначенные задачи сохранены: тесты не исправляются, невыгодные случаи после просмотра результатов не удаляются. Поэтому отсутствие потолка относится к \emph{наблюдаемому выполнению тестов}; оно само по себе не доказывает внутреннюю сложность задач или правильность программ относительно условия. Перед подтверждением нужны аудит валидности без просмотра исходов моделей и независимый бенчмарк.

В опубликованных доказательствах первоначальная остановленная попытка отделена от заменяющей, её токены учтены. Арифметические пробы и исключённая проверка новой среды также не входят в экспериментальную матрицу. Изменение среды после инфраструктурного сбоя на development не создаёт независимую репликацию. Большие отложенные выборки, повторные генерации, чувствительность при равном бюджете и настоящие внешние baseline остаются необходимыми для закрытия существенных эмпирических замечаний.
